% Auriga theme
% https://github.com/anishathalye/auriga

\documentclass[14pt,aspectratio=169]{beamer}
\usepackage{pgfpages}
\usepackage{fancyvrb}
\usepackage{tikz}
\usepackage{pgfplots}
\usepackage[slovene]{babel}
\usepackage{amsmath}
\usepackage{mathdots}

\newcommand{\R}{\mathbb R}
\newcommand{\N}{\mathbb N}
\newcommand{\Z}{\mathbb Z}
\newcommand{\loss}{\mathcal L}
\newcommand{\fun}{\operatorname{fun}}
\newcommand{\funnn}[1]{\fun([#1], [#1])}
\newcommand{\Loss}[1]{\mathcal L _\text{#1}}
\newcommand{\C}{\mathbb C}
\newcommand{\Q}{\mathbb Q}
\newcommand{\mon}{(S\cup S ^{-1})^*}
\newcommand{\s}{\hat \rho (s)}

\theoremstyle{definition}
\newtheorem{aksiom}{Aksiom}
\newtheorem{definicija}{Definicija}[section]
\newtheorem{opomba}[definicija]{Opomba}
\newtheorem{primer*}[definicija]{Primer}
\theoremstyle{plain}
\newtheorem{lema}[definicija]{Lema}
\newtheorem{izrek}[definicija]{Izrek}
\newtheorem{trditev}[definicija]{Trditev}
\newtheorem{posledica}[definicija]{Posledica}

%NOTES
%\setbeamertemplate{note page}[plain]
%\setbeameroption{show notes on second screen=right}


\usetheme{auriga}
\usecolortheme{auriga}

% define some colors for a consistent theme across slides
\definecolor{red}{RGB}{181, 23, 0}
\definecolor{blue}{RGB}{0, 118, 186}
\definecolor{gray}{RGB}{146, 146, 146}

\definecolor{labsixorange}{RGB}{243, 111, 33}
\definecolor{aoblue}{RGB}{8, 96, 204}

\title{Iskanje morfizmov z gradientnim spustom}

\author{{Urh Primožič} \inst{1} \and Ljupčo Todorovski \inst{2} \and Urban Jezernik \inst{2}}

\institute[shortinst]{\inst{1} mentor \samelineand \inst{2} somentor}

\begin{document}

{
  % rather than use the frame options [noframenumbering,plain], we make the
  % color match, so that the indicated page numbers match PDF page numbers
  \setbeamercolor{page number in head/foot}{fg=white}
  \begin{frame}
    \titlepage
  \end{frame}
}

\begin{frame}{\color{aoblue}Diskretni objekti so ničle funkcionalov}
\note{
Razloži osnovno idejo. Matematične objekte vpeno v evklidkski prostor in jih enačimo z množico ničel 
nekega funkcionala. Potem lahko te objekte iščemo preko optimizacije.
}
    \begin{itemize}
        \item iščemo objekte $Z  \subset Y^X$
        \item vpetje  $Y^X \hookrightarrow  \R^N$
        \item $Z = \loss^{-1}(0)$ za $\loss \colon \R^N \to [0,\infty)$
        \item $\operatorname{argmin}\loss \in Z$
    \end{itemize}
\end{frame}
\begin{frame}{\color{aoblue}Modeliranje upodobitev}
  \note{
    Poljubna preslikava je upodobitev, če je izguba enaka 0.
  }

  \begin{align*}
    \rho \colon G &\to \mathbb{C}^{n \times n} \\
    g &\mapsto
    \begin{bmatrix}
      g_{1,1} & \cdots & g_{1,n} \\
      \vdots & & \vdots \\
      g_{n,1} & \cdots & g_{n,n}
    \end{bmatrix} = \phi_g
  \end{align*}
  je upodobitev, če je 
  $$
  \left\|
    \rho(gh) - \rho(g)\rho(h)
  \right\| = 0
  $$
  za vse $g, h \in G$.
\end{frame}
\begin{frame}{\color{aoblue}Modeliranje upodobitev}
  \note{
    Začnemo s poljubno preslikavo, in jo postopoma spreminjamo, da bo bolj podobna upodobitvi, kar 
    je ekvivalentno temu, da je izguba čim manjša.
  }
Poljubna preslikava 
  \begin{align*}
   \hat {\rho} \colon G &\to \mathbb{C}^{n \times n} \\
    g &\mapsto
    \begin{bmatrix}
      g_{1,1} & \cdots & g_{1,n} \\
      \vdots & & \vdots \\
      g_{n,1} & \cdots & g_{n,n}
    \end{bmatrix} = \phi_g
  \end{align*}
  z izgubo 
  $$
  \sum_{g, h \in G} \left\|
    \hat{\rho}(gh) - \hat{\rho}(g)\hat{\rho}(h)
  \right\|
  $$
\end{frame}
\begin{frame}{Generatorji}
  \note{
    V praksi hočemo čim manj parametrov, zato se omejimo le na generatorji. 
  }

  Grupa s prezentacijo $G = <S | R>$. Na $S$:
  $$
  \hat{\rho}(s) =
  \begin{bmatrix}
    s_{1,1} & \cdots & s_{1,n} \\
    \vdots & & \vdots \\
    s_{n,1} & \cdots & s_{n,n}
  \end{bmatrix} = \phi_s
  $$ 
  drugje implicitno.
  % \vfill
  % \color{gray}{Model $\hat \rho$ has parameters $\phi = \{\phi_s
  % \mid s \in S\} $. }
\end{frame}
\begin{frame}{Funkcije izgube}
        \begin{align*}
            \Loss{rel}(\hat\rho) &= \frac{1}{|R|} \sum \limits_{r \in R}
    \|\hat \rho(r) - I  \|^2 
     &\text{ \small izguba na relacijah}\\ 
    \Loss{irr}(\rho) &= (|\chi_\rho| - 1)^2 
     &\text{ \small nerazcepnost}\\
     \Loss{unitary} &= \frac{1}{|S|} \sum \limits_{s \in S}
    \|\hat\rho(s) \hat\rho(s)^H - I\|^2
     &\text{ \small unitarnost}
    \end{align*}
    
$$
\loss = \Loss{rel} + \Loss{irr} + \Loss{unitary}
$$
\end{frame}

\begin{frame}{\color{aoblue}Gradientni spust}
\note{

}
Optimizacija z gradientnim spustom:
    $$
    \phi_{i+1} = \phi_i - \eta \nabla \loss(\phi_i) 
    $$
    Enačba gradientnega toka:
    $$
    \frac{d\phi}{dt} = - \nabla \loss(\phi), \quad \phi(0) = \phi_0
    $$
\end{frame}
\begin{frame}
    \begin{center}
        \includegraphics[height=0.7\paperheight ]{gradientni_spust_primer.pdf}
$$
f(x) = \sin(x)\sqrt{x+2}
$$
    \end{center}
\end{frame}

\begin{frame}{Primer}
    Realen model $ \hat  \rho \colon S_3 = \langle (12), (13) \rangle \to \R$. 
    \begin{center}
        \includegraphics[height=0.7\paperheight ]{img/demo_S3_ena_trajektorija.pdf}
\end{center}
\end{frame}
\begin{frame}{Primer}
    Realen model $ \hat  \rho \colon S_3 = \langle (12), (13) \rangle \to \R$. 
    \begin{center}
        \includegraphics[height=0.7\paperheight ]{img/demo_trajektorije_S3_to_R.pdf}
\end{center}
\end{frame}
\begin{frame}{\color{aoblue}Obstoj rešitve}
  \begin{izrek}
    Če je $\loss(\hat \rho) \leq O(\epsilon)$, obstaja
    \emph{upodobitev} $\rho$, da velja
    $$
    sup_{g \in G} \left\|
    \hat \rho(g) - \rho(g)
    \right\| \leq \epsilon.
    $$
  \end{izrek}
\end{frame}
\begin{frame}{$\epsilon$-upodobitve}
  \begin{izrek}[Kazhdan]
    Če za $\mu \colon G \to U(\mathcal H)$ iz amenabilne grupe velja
    $$
    \sup_{g,h \in G} \|\mu(gh) - \mu(g)\mu(h)\| < \epsilon \leq \frac{1}{100},
    $$
    obstaja unitarna upodobitev $\pi$ grupe $G$, da velja
    $$
    \sup_{g \in G} \|\pi(g) - \mu(g) \| \leq 2\epsilon.
    $$
  \end{izrek}
\end{frame}
\begin{frame}{kvaziupodobitve}
  \begin{definicija}
    $\hat \rho \colon \langle S \mid R \rangle \to \C^{n \times n}$ je
    \emph{kvaziupodobitev}, če velja
    $$
    sup_{r \in R} \left\| \phi(r) - I  \right\| < \epsilon.
    $$
  \end{definicija}
  \begin{izrek}
    Naj bo $G$ končna grupa. $\hat \rho \colon G \to \C^{n \times n}$
    je  $\epsilon$-upodobitev natanko tedaj, ko je kvaziupodobitev.
  \end{izrek}
\end{frame}
\begin{frame}{\color{aoblue}Ciklične grupe}
    \note{
        $C_n$ ima le en generator $z$, vse nerazcepne upodobitve so enodimenzionalne in kompleksne. 
        Torej iščeš $\hat\rho(z)=x+iy \in  \C$.
    }
    Nerazcepne
upodobitve $C_n=\langle z \mid z^n = 1 \rangle$ so
$$
\rho_k(z) = e^{\frac{2 \pi i k}{n}}.
$$
Funkcije izgube:
\begin{align*}
      \Loss{rel} &= |(x + iy)^n -1|^2\\
  \Loss{irr} &= (\frac{1}{n} \sum_{j=0}^{n-1} |(x + iy)^j|^2 -1)^2 \\
  \Loss{unitary} &= | x^2 + y^2 -1|^2.
\end{align*}
\end{frame}
\begin{frame}{Graf izgube}
    \centering
    \includegraphics[height=0.72\paperheight]{img/C3 landscape.pdf}
\end{frame}
\begin{frame}{Trajektorije gradientnega spusta}
    \centering
    \includegraphics[height=0.72\paperheight]{img/C3 trajektorije.pdf}
\end{frame}
\begin{frame}{\color{aoblue}Diedrske grupe}
    $D_n = \langle r, s \mid r^n, s^2, (rs)^2 \rangle$
\begin{align*}
  \chi_{\varepsilon, \delta} \colon \quad
  s &\mapsto \varepsilon, \\
  r &\mapsto \delta, \quad  \varepsilon \in \{-1, 1\}, \delta \in \{1, -1^{n + 1}\}
\end{align*}
\begin{align*}
  \rho_k \colon  s &\mapsto
  \begin{bmatrix}
    0 & 1 \\
    1 & 0
  \end{bmatrix}\\
  r &\mapsto
  \begin{bmatrix}
    \cos\left(\frac{2\pi k}{n}\right) & -\sin\left(\frac{2\pi
    k}{n}\right) \\
    \sin\left(\frac{2\pi k}{n}\right) & \cos\left(\frac{2\pi k}{n}\right)
  \end{bmatrix}
\end{align*}
\end{frame}
\begin{frame}{Enodimenzionalne upodobitve}
    \begin{center}
        \includegraphics[height=0.8\paperheight ]{img/Dn_1dim_n1.png}
    \end{center}
\end{frame}
\begin{frame}{Enodimenzionalne upodobitve}
    \begin{center}
        \includegraphics[height=0.8\paperheight ]{img/Dn_1dim_n2.png}
    \end{center}
\end{frame}
\begin{frame}{Enodimenzionalne upodobitve}
    \begin{center}
        \includegraphics[height=0.8\paperheight ]{img/Dn_1dim_n7.png}
    \end{center}
\end{frame}
\begin{frame}{Enodimenzionalne upodobitve}
    \begin{center}
        \includegraphics[height=0.8\paperheight ]{img/Dn_1dim_n8.png}
    \end{center}
\end{frame}
\begin{frame}{Dvodimenzionalne upodobitve}
    \begin{center}
        \includegraphics[height=0.8\paperheight ]{img/prob_of_convergence_D_3.png}
    \end{center}
\end{frame}
\begin{frame}{Dvodimenzionalne upodobitve}
    \begin{center}
        \includegraphics[height=0.8\paperheight ]{img/prob_of_convergence_D_4.png}
    \end{center}
\end{frame}
\begin{frame}{Dvodimenzionalne upodobitve}
    \begin{center}
        \includegraphics[height=0.8\paperheight ]{img/prob_of_convergence_D_11.png}
    \end{center}
\end{frame}
\begin{frame}{Dvodimenzionalne upodobitve}
    \begin{center}
        \includegraphics[height=0.8\paperheight ]{img/prob_of_convergence_D_12.png}
    \end{center}
\end{frame}
\begin{frame}
    \centering
    \includegraphics[height=\paperheight ]{img/plot_grid_n3_-1_1_r5000.pdf}
\end{frame}
\begin{frame}
    \centering
    \includegraphics[height=\paperheight ]{img/cut_plot_grid_n3_0.5_1.5_r3000.pdf.png}
\end{frame}
\begin{frame}
    \centering
    \includegraphics[height=\paperheight ]{img/cut_plot_grid_n3_-1.5_-0.5_r2000.pdf.png}
\end{frame}
\begin{frame}
    \centering
    \includegraphics[height=\paperheight ]{img/plot_grid_n4_-1_1_r2000.pdf}
\end{frame}
\begin{frame}
    \centering
    \includegraphics[height=\paperheight ]{img/plot_grid_n4_0.5_1.5_r2000.pdf.png}
\end{frame}
\begin{frame}{\color{aoblue}Delovanja}
  \note{
    Lets adapt the approach to group actions.

    Show some motivation on the board.

    Discrete space is not suitable for GD - work with probabilities!
    Explain how we get to the matrices, explain that this is not the only way
  }
  Elementi grupe $G=<S|R>$ so slučajne preslikave $g \colon [n] \to [n]$.
  \begin{align*}
    \hat \rho \colon S &\to \mathcal{P}(\operatorname{fun}([n],[n])) \\
    s &\mapsto P_s =
    \begin{bmatrix}
      P(s(1) = 1) & P(s(1) = 2) & \dotsm & P(s(1) = n) \\
      P(s(2) = 1) & P(s(2) = 2) & \dotsm & P(s(2) = n) \\
      \vdots & \vdots & & \vdots \\
      P(s(n) = 1) & P(s(n) = 2) & \dotsm & P(s(n) = n) \\
    \end{bmatrix}
  \end{align*}
  Za $f \in \operatorname{fun}([n],[n])$ je
  $$P(s = f) = \prod \limits_{i=1}^n P(s(i) = f(i)) = \prod
  \limits_{i=1}^n s_{i, f(i)}$$
\end{frame}
\begin{frame}{Parametrizacija stohastičnih matrik}
  \begin{itemize}
    \item vektor v verjetnostno
      porazdelitev:$\operatorname{softmax}(v)_i =
      \frac{e^{v_i}}{\sum_{j=1}^n e^{v_j}}$
      % pretvori
      %vektor $v$ v  verjetnostno porazdelitev.
    \item parametrizacija stohastičnih matrik:
      $$
      \begin{bmatrix}
        a_1^T \\
        \vdots \\
        a_n^T
      \end{bmatrix} \mapsto
      \begin{bmatrix}
        \operatorname{softmax}(a_1^T) \\
        \vdots \\
        \operatorname{softmax}(a_n^T)
      \end{bmatrix}
      $$
  \end{itemize}
\end{frame}
\begin{frame}{\color{aoblue} Izguba na relacijah}
  Za relacijo $r \in R$ hočemo $P(r=\operatorname{id}) = 1$
  $$
  0 = \log(P(r=\operatorname{id})) = \log(\prod_{i=1}^n P(r(i)=i)) =
  \operatorname{tr}(\log(P_r))
  $$
  $$
  \mathcal{L}_{\text{rel}}  = - \frac{1}{|R|} \sum_{r \in R}
  \operatorname{tr}(\log(P_r))
  $$
\end{frame}
\begin{frame}{$C_6 = \langle z \mid z^6 \rangle$ deluje na $[6]$}
    \begin{columns}
        \begin{column}{0.5\textwidth}
            %\begin{figure}
                \centering
                \includegraphics[width=1.2\linewidth]{img/C6_on_6_sample2_initial.pdf}
          %  \end{figure}
        \end{column}
        \begin{column}{0.5\textwidth}
            %\begin{figure}
                \centering
                \includegraphics[width=1.2\linewidth]{img/C6_on_6_sample2_final.pdf}
        % \end{figure}
        \end{column}
    \end{columns}
     \centering 
     $\hat \rho(z)$ konvergira do $ (1 6 2  3   5 4 )$
\end{frame}
\begin{frame}{$D_5 = \langle r, s \mid r^n, s^2, (rs)^2 \rangle$ deluje na $[5]$}
\begin{figure}
    \centering
    \includegraphics[height=0.7\textheight]{img/D5_on_5_sample7_initial.pdf}
    % \caption{Caption}%
\end{figure}
\end{frame}
\begin{frame}{$D_5 = \langle r, s \mid r^n, s^2, (rs)^2 \rangle$ deluje na $[5]$}
\begin{figure}
    \centering
    \includegraphics[height=0.7\textheight]{img/D5_on_5_sample7_final.pdf}
    $r$ konvergira k  $(13524)$, $s$  konvergira k  $(14)(23)$
\end{figure}
\end{frame}
\begin{frame}{Prehod na delovanje}
    \begin{itemize}
        \item $\overline{P_s}$ - $1$ na največjem elementu v vrstici, $0$ drugje
        \item pri majhni izgubi je $s \mapsto \overline{P_s}$ delovanje
    \end{itemize}
\end{frame}
\begin{frame}{\color{aoblue} Izomorfizmi grafov}
\begin{itemize}
    \item  $M_1 = \begin{bmatrix}
        m_{i,j}^{(1)}
    \end{bmatrix}$, $M_2 = \begin{bmatrix}
        m_{i,j}^{(2)}
    \end{bmatrix}$ matriki sosednosti 

    \item slučajna preslikava $f \colon [n] \to [n] \sim P_f = \begin{bmatrix}P(f(i)=j)
    \end{bmatrix}$

    
    \item za $i \sim_1 j$ hočemo $P(f(i) \sim_2 f(j) = 1$ 
\end{itemize}    

\end{frame}
\begin{frame}{\color{aoblue}Izguba na relacijah}
 
$$
P(f(i) \sim_2 f(j)) = \sum_{k= 1}^n \sum_{h = 1}^n f_{i, k} f_{j, h} m_{k,h}^{(2)} 
% = f_j^T M_1 f_i
= (P M_2 P^T)_{i, j}
$$  
$$
 \loss_\sim = \|M_1 - P M_2 P^T\|.
$$
\end{frame}
\begin{frame}{\color{aoblue}Iščemo bijekcije}
  \note{Similarly to the unitarity constraint, we can constraint bijectivity}
  $$
  P(f \in S_n )= \sum_{\sigma \in S_n} P(f = \sigma) = \sum_{\sigma
  \in S_n}   \prod_{i=1}^n f_{i, \sigma(i)}
  = \text{Perm} (P_f).
  $$
  $P(f \in S_n ) = 1 \iff P_f \text{ permutacijska }\iff P_f \in U_n$
  \vfill
  $$
  \mathcal{L}_{\text{bijective}} = \frac{1}{|S|} \sum_{f \in
  S}||P_fP_f^T - I||_F^2  = \mathcal{L}_{\text{unitary}} \circ
  \operatorname{softmax}
  $$
\end{frame}
\begin{frame}{Primer}
    \begin{columns}
        \begin{column}{0.5\textwidth}
            \begin{figure}
                \centering
                \includegraphics[, width=0.7\textwidth]{img/graf_primer.png}
              %  \caption{Caption}
                \label{fig:enter-label}
            \end{figure}
        \end{column}
        \begin{column}{0.5\textwidth}
            \begin{figure}
                \centering
                \includegraphics[height=0.6\textheight]{img/graf_primer_final.pdf}
              %  \caption{Caption}
                
            \end{figure}
            Model konvergira do $(13)(24)$
        \end{column}
    \end{columns}
\end{frame}
\begin{frame}{\color{aoblue}Tabela inverzij}
  \begin{definicija}[Tabela inverzij]
    $$\operatorname{TI}(\sigma) =
    (a_1,
      a_2, \dotsc,
    a_n), \quad a_k = |\{ j > k
      \mid \sigma(j) < \sigma(k)
    \}|$$
  \end{definicija}
  \begin{izrek}
    $\operatorname{TI} \colon S_n \to \prod_{i=1}^n \{0, 1, \dotsc, n-i\}$ je bijekcija.
  \end{izrek}
  \end{frame}
\begin{frame}{Porazdelitev nad tabelo inverzij}
    \centering
    $$
    P_d =
    \begin{bmatrix}
      p_{1,1} & p_{1,2} & \dotsm & p_{2, n-1} & p_{1,n} \\
      p_{2,1} & p_{2,2} & \dotsm & p_{2, n-1} & 0 \\
      \vdots & \vdots & \iddots & \vdots & \vdots    \\
      p_{n-1,1} & p_{n-1,2}  & \dotsm & 0  & 0\\
      p_{n,1} & 0  & \dotsm & 0  & 0\\
    \end{bmatrix}
    $$
    $$P(a_k = i) = P(k \text{ postavimo na } i\text{-to prazno
mesto}) = p_{k, i}$$
$$P(\operatorname{TI}(\sigma) = (x_1, x_2, \dotsc, x_n)) = \prod_{i=1}^n
P(a_k = x_k)$$
\end{frame}
\begin{frame}{Funkcija izgube}
    \begin{itemize}
        \item $\Loss{TI} = \sum_{i, j} ( (M_1)_{i,j} - P(f(i) \sim f(j)) )^2$
        \item preparametrizacija $P_d = M(\phi)(x_0)$ 
    \end{itemize}
\end{frame}
\begin{frame}{Rezultati}
    \centering
    \includegraphics[height=0.8\paperheight]{img/vanilla_vs_TI_loss.pdf}
\end{frame}
\begin{frame}{Rezultati}
    \centering
    \includegraphics[height=0.8\paperheight]{img/average_curve_full.pdf.pdf}
\end{frame}
\begin{frame}{Opombe}
    \begin{itemize}
        \item $\epsilon$-upodobitve in kvaziupodobitve
        \item softmax vs Sinkhorn
        \item simetrije
        \item meta učenje in boljše arhitekture modelov
    \end{itemize}
\end{frame}
\end{document}
